\newcommand{\resReviewScore}{0}
\newcommand{\resReviewVerbal}{brak oceny, trwa pierwsza iteracja}
\newcommand{\resReviewHistory}{brak}


\begin{frame}[fragile]{Recenzent AI ocenia talię według trzech źródeł kryteriów tym samym promptem w~każdej iteracji}
  \begin{itemize}
    \item Agent presentation-reviewer (opus, effort high) czyta PDF, źródła .tex, strony PNG i~wyniki poleceń sprawdzających. Jest tylko do odczytu na mocy instrukcji.
    \item Kryteria pochodzą z~zasad zaliczenia, instrukcji projektu 2 i~przewodnika stylu prezentacji.
    \item Rdzeń promptu poniżej jest ten sam w~każdej iteracji. Zmienia się tylko akapit ze ścieżkami plików i~numerem iteracji.
  \end{itemize}
\begin{Verbatim}[commandchars=\\\{\},fontsize=\footnotesize]
[...]
Review slides/preliminary/main.pdf and its .tex sources against three sources
of criteria. (1) materials/zasady_zaliczenia_projektu_v1.pdf, the mandatory
conditions and slide groups. (2) materials/Projekt_2.pdf, section 4 on page 3,
the content of the preliminary deck. (3) docs/presentation-style.md,
assertion-evidence titles, text budget, palette, banned elements. Report
findings with severity and a suggested fix. End with a verbal grade and a
score 1-10.
[...]
\end{Verbatim}
\end{frame}

\begin{frame}{Recenzent AI ocenił prezentację wstępną na \resReviewScore{} / 10}
  „\resReviewVerbal”
  \vfill
  {\color{mnMuted}\footnotesize Przebieg iteracji: \resReviewHistory\par}
\end{frame}
